% Lab 1 — PrimeTime Flow（原书 PDF p.3–18）
\bichapter{Lab 1：PrimeTime 流程}{Lab 1: PrimeTime Flow}

\begin{objectiveBox}
完成本实验后，你应能够：
\begin{itemize}
  \item 恢复（restore）先前保存的 PrimeTime 会话（session）；
  \item \textbf{（可选）} 使用若干有助于交互式效率的 PrimeTime 命令，并学会如何查找命令、变量与设计库相关信息；
  \item 验证已恢复的 save-session；
  \item 练习推荐的 PrimeTime 流程（flow）；
  \item 解读建立时间（setup）与保持时间（hold）时序报告中的关键组成部分。
\end{itemize}
\end{objectiveBox}

\begin{durationBox}
约 \textbf{45} 分钟。
\end{durationBox}

\section{概述}
\label{lab1:overview}

\noindent\textit{Overview}

本实验大致步骤如下：
\begin{enumerate}
  \item 恢复一个 PrimeTime 会话；
  \item 学习若干有用命令（\textbf{可选任务}）；
  \item 验证已恢复的保存会话，并练习推荐流程；
  \item 解读建立时间与保持时间时序报告。
\end{enumerate}

\subsection{相关文件与目录}
\label{lab1:files}

\noindent\textit{Relevant Files and Directories}

本实验全部文件位于主目录下的 \cmd{lab1_flow} 目录中。

\begin{tabularx}{\textwidth}{@{}l X@{}}
\toprule
\textbf{路径 / 文件} & \textbf{说明} \\
\midrule
\cmd{lab1_flow/} & 当前工作目录 \\
\cmd{common_setup.tcl} & 多工具共享的 setup 文件 \\
\cmd{pt_setup.tcl} & PrimeTime 专用 setup 文件 \\
\cmd{pt_scripts/} & 运行脚本目录 \\
\quad\cmd{pt.tcl} & 运行脚本（run file） \\
\cmd{.synopsys_pt.setup} & 自动读取的 PT setup 文件 \\
\cmd{orca_savesession} & 已保存会话目录 \\
\cmd{RUN.tcl} & 用于重建 \cmd{orca_savesession} 的运行脚本 \\
\bottomrule
\end{tabularx}


\section{操作说明}
\label{lab1:instructions}

\noindent\textit{Instructions}

目标：通过验证已恢复的 save session、练习推荐流程，并分析 setup/hold 报告，熟悉 PrimeTime。

\subsection{任务 1：恢复 PrimeTime 会话}
\label{lab1:task1}

\noindent\textit{Task 1. Restore a PrimeTime Session}

调用先前保存的 PrimeTime 会话以执行静态时序分析（STA, Static Timing Analysis）。

\begin{enumerate}
  \item 从实验 Unix 目录 \cmd{lab1_flow} 启动 PrimeTime：
\begin{lstlisting}
unix% cd lab1_flow
unix% pt_shell
\end{lstlisting}

  \item 恢复先前保存的 PrimeTime 会话。该步骤将读入设计网表（netlist）、库（libraries）与约束（constraints）。此后设计即可进行分析。

\begin{noteBox}
\begin{itemize}
  \item 下文的 \cmd{orca_savesession} 是一个 Unix 目录。
  \item 如需重建该会话，可使用：\\
        \cmd{pt_shell -f RUN.tcl | tee -i run.log}
  \item 执行 \cmd{RUN.tcl} 时出现的任何 \textbf{PARA-124} 错误，在本课程实验中可安全忽略。
  \item PrimeTime 支持命令、选项、变量与文件名补全：键入少量字符后按 \textbf{Tab} 键即可。
\end{itemize}
\end{noteBox}

\begin{lstlisting}
pt_shell> restore_session orca_savesession
\end{lstlisting}

  \item 生成覆盖率分析报告：
\begin{lstlisting}
pt_shell> report_analysis_coverage
\end{lstlisting}

\begin{questionBox}
\textbf{问题 1.}\ 当前分析的设计名称是什么？\\[0.35em]
\textbf{答：} 设计名称为 \textbf{ORCA}。

\textbf{问题 2.}\ ORCA 有多少个 setup 违例与 hold 违例？\\[0.35em]
\textbf{答：} 共有 \textbf{23} 个 setup 违例与 \textbf{53} 个 hold 违例。
\end{questionBox}

  \item 生成全局时序报告：
\begin{lstlisting}
pt_shell> report_global_timing
\end{lstlisting}

\begin{questionBox}
\textbf{问题 3.}\ 其中有多少个 reg-reg 类型的 setup 与 hold 违例？\\[0.35em]
\textbf{答：} 类型为 reg-reg 的违例为：\textbf{9} 个 setup、\textbf{53} 个 hold。
\end{questionBox}
\end{enumerate}

\subsection{任务 2：［可选］探索有用命令}
\label{lab1:task2}

\noindent\textit{Task 2. [OPTIONAL TASK] Explore Helpful Commands}

\begin{enumerate}
  \item 执行以下三条历史快捷命令：
\begin{lstlisting}
pt_shell> history
pt_shell> !!
pt_shell> !2
\end{lstlisting}

\begin{questionBox}
\textbf{问题 4.}\ 说明上面最后两条历史命令的区别。\\[0.35em]
\textbf{答：} 命令 \cmd{!2} 重复历史中第 2 条已执行命令；命令 \cmd{!!} 重复历史中最近一条已执行命令。
\end{questionBox}

  \item 也可用上下箭头在历史事件列表中滚动，作为上一步的替代方式。\\
        执行以下命令查看默认 emacs 编辑模式下的全部键绑定：
\begin{lstlisting}
pt_shell> list_key_bindings
\end{lstlisting}

  \item 体验 page mode 别名：执行下列命令生成会滚出屏幕的长报告：
\begin{lstlisting}
pt_shell> report_timing -group [get_path_group *]
\end{lstlisting}

  \item 打开 page mode：
\begin{lstlisting}
pt_shell> page_on
pt_shell> !rep
\end{lstlisting}

  \item 用空格键翻页阅读长报告；在 page mode 下键入 \texttt{q} 可退出长报告。
        若要关闭 page mode，使用别名 \cmd{page_off}。

  \item 使用 Tcl 过程 \cmd{view}，把时序报告送到单独窗口：
\begin{lstlisting}
pt_shell> view report_timing -group [get_path_group *]
\end{lstlisting}

  \item 查找用于恢复 PrimeTime 会话的命令，并显示该命令的帮助信息：
\begin{lstlisting}
pt_shell> help restore*
pt_shell> man restore_session
pt_shell> restore_session -help
\end{lstlisting}

\begin{noteBox}
以下是显示语法帮助的另一种方式：
\begin{lstlisting}
pt_shell> help -v restore_session
\end{lstlisting}
\end{noteBox}

\begin{questionBox}
\textbf{问题 5.}\ 根据上一条命令，\cmd{restore_session} 是否接受开关（switches）？\\[0.35em]
\textbf{答：} 是。它接受会话名等参数/开关，从而可指定恢复众多已保存会话中的哪一个。
\end{questionBox}

  \item PrimeTime 中的时间单位由主工艺库（main technology library）决定。
        要查找 ORCA 的时间单位，先列出内存中的全部库。

\begin{noteBox}
报告中库名前的 \texttt{*} 表示主库（main library）。
\end{noteBox}

\begin{lstlisting}
pt_shell> list_lib
\end{lstlisting}

  \item 对主库生成报告以查看时间单位。

\begin{noteBox}
建议复制粘贴库名以免拼写错误。时间单位位于报告最顶部。
\end{noteBox}

\begin{lstlisting}
pt_shell> report_lib cb13fs120_tsmc_max
\end{lstlisting}

\begin{questionBox}
\textbf{问题 6.}\ ORCA 设计的时序报告（以及所有其他报告）使用的时间单位是什么？\\[0.35em]
\textbf{答：} \textbf{1\,ns}。
\end{questionBox}

\begin{noteBox}
别忘了在 page mode 下用 \texttt{q} 退出长报告，无需读完整份报告即可回到 \cmd{pt_shell} 提示符。
\end{noteBox}

  \item 显示当前设计使用的单位：
\begin{lstlisting}
pt_shell> report_units
\end{lstlisting}
\end{enumerate}

\subsection{任务 3：验证已有的 PrimeTime 会话}
\label{lab1:task3}

\noindent\textit{Task 3. Validate an Existing PrimeTime Session}

本任务验证已读入 PrimeTime 的输入：当前设计与库、反标（back-annotation）与约束。

\begin{enumerate}
  \item 确认当前设计为顶层模块 \textbf{ORCA}：
\begin{lstlisting}
pt_shell> current_design
\end{lstlisting}

  \item 将 Unix 路径上的库与已读入 PrimeTime 的库进行对照：
\begin{lstlisting}
pt_shell> printvar search_path
pt_shell> printvar link_path
pt_shell> list_libraries
\end{lstlisting}

\begin{questionBox}
\textbf{问题 7.}\ \cmd{link_path} 中的 4 个库是否都已成功读入 PrimeTime？\\[0.35em]
\textbf{答：} 是——用 \cmd{link_library} 指定的四个库与 \cmd{list_libraries} 显示的一致。

\textbf{问题 8.}\ 哪一个库定义了时间单位、工作条件（operating conditions）以及其他延迟计算相关信息的默认值？\\[0.35em]
\textbf{答：} \cmd{cb13fs120_tsmc_max}——\cmd{list_libraries} 输出中库名前的星号表示“主库”。

\textbf{问题 9.}\ 使用的时间单位是什么？\\[0.35em]
\textbf{答：}\\
\begin{lstlisting}
pt_shell> report_lib cb13fs120_tsmc_max
...
Time Unit : 1ns

pt_shell> report_units
...
Time_unit : 1e-09 Second
\end{lstlisting}
\end{questionBox}

  \item 验证网线是否已完整反标：
\begin{lstlisting}
pt_shell> report_annotated_parasitics
\end{lstlisting}

\begin{questionBox}
\textbf{问题 10.}\ 是否存在未反标的网？\\[0.35em]
\textbf{答：} 是：存在一些内部无驱动网（driverless nets），以及一些边界 pin-to-pin 网未反标。

\textbf{问题 11.}\ 作为调试缺失网的“下一步”，\cmd{report_annotated_parasitics} 的哪个选项比较合适？\\[0.35em]
\textbf{答：} \cmd{report_annotated_parasitics -help} 可列出选项；
\opt{-list_not_annotated} 可列出网名；
若关注 pin-to-pin 网，可使用 \opt{-pin_to_pin_nets}。
\end{questionBox}

  \item 验证设计是否已完整约束：
\begin{lstlisting}
pt_shell> check_timing
\end{lstlisting}

\begin{questionBox}
\textbf{问题 12.}\ 作为调试缺失约束的“下一步”，\cmd{check_timing} 的哪个选项比较合适？\\[0.35em]
\textbf{答：} \cmd{check_timing -help} 列出选项；\opt{-verbose} 可列出未约束端点（unconstrained endpoints）。
\end{questionBox}

  \item 验证单元中的检查是否被充分行使（exercised）；并思考可能原因：
\begin{lstlisting}
pt_shell> report_analysis_coverage
pt_shell> report_case_analysis
\end{lstlisting}

\begin{questionBox}
\textbf{问题 13.}\ 许多时序检查处于 untested 状态——这是否合理？\\[0.35em]
\textbf{答：} 合理——\cmd{test_mode}、\cmd{scan_en} 与 \cmd{power_save} 被约束为关；
对额外模式做 STA 可能使更多检查变为 tested。
\end{questionBox}

  \item 退出 PrimeTime：
\begin{lstlisting}
pt_shell> quit
\end{lstlisting}
\end{enumerate}

\subsection{任务 4：执行运行脚本并分析运行结果}
\label{lab1:task4}

\noindent\textit{Task 4. Execute the Run Script and Analyze the Run}

\begin{enumerate}
  \item 执行运行脚本，并将结果记录到日志文件 \cmd{run.log}：
\begin{lstlisting}
UNIX> pt_shell -f ./pt_scripts/pt.tcl | tee -i run.log
\end{lstlisting}

\begin{questionBox}
\textbf{问题 14.}\ 执行运行脚本过程中是否出现错误？\\[0.35em]
\textbf{答：} 没有。该运行脚本的设计使得一旦出错会在中途终止；若脚本完整跑完，则执行过程中未发生错误。
此外，日志末尾由 \cmd{quit} 给出的 Diagnostics Summary 也显示本次运行无错误。
\end{questionBox}

  \item 若有错误，须先处理错误，再进入下一步。

  \item 评估日志文件：用文本编辑器打开日志，搜索 \cmd{update_timing} 相关消息（\textbf{UITE-214}）、
        \cmd{print_message_info} 输出，以及 \cmd{quit} 输出。
        然后在 profile 目录中查看文件 \cmd{tcl_profile_sorted_by_cpu_time}。

\begin{questionBox}
\textbf{问题 15.}\ 哪一步占用的 CPU 时间最多？\\[0.35em]
\textbf{答：} 对本实验而言，\textbf{source 约束}占用的 CPU 时间最多。

\textbf{问题 16.}\ 哪些命令引发了 UITE-214 消息？\\[0.35em]
\textbf{答：} \cmd{update_timing -full}，以及调用宏 \cmd{PLL_SHIFT}（其内部触发了 \cmd{update_timing}）。

\textbf{问题 17.}\ 是否有部分 timing update 可以避免？\\[0.35em]
\textbf{答：} 可以。在 \cmd{set_propagated_clocks} 前后各有一次 \cmd{update_timing}——
之前那一次是不必要的。一次 \cmd{update_timing} 位于 \cmd{pt.tcl}，另一次位于约束文件
\cmd{orca_pt_other.tcl}（由 \cmd{orca_pt_constraints} source，再由 \cmd{pt.tcl} source）。

\textbf{问题 18.}\ 为什么在通读日志之前，从 \cmd{quit} 命令的输出入手可能是个好起点？\\[0.35em]
\textbf{答：} 它给出潜在问题点的高层摘要：警告/信息消息、timing update，以及性能统计。
\end{questionBox}
\end{enumerate}

\subsection{任务 5：分析 STA 报告}
\label{lab1:task5}

\noindent\textit{Task 5. Analyze STA Reports}

为时钟 \cmd{SYS_CLK} 生成并解读 setup 与 hold 两份 STA 报告。

\begin{enumerate}
  \item 启动 PrimeTime，并恢复上一任务中保存的会话：
\begin{lstlisting}
unix% pt_shell
pt_shell> restore_session my_savesession
\end{lstlisting}

  \item 显示 ORCA 中的时钟：
\begin{lstlisting}
pt_shell> report_clock
\end{lstlisting}

\begin{questionBox}
\textbf{问题 19.}\ ORCA 中有多少个时钟？\\[0.35em]
\textbf{答：} ORCA 中有 \textbf{6} 个时钟。
\end{questionBox}

  \item 为时钟 \cmd{SYS_CLK} 生成一份“简短（short）”的 setup 时序报告。
        使用命令行补全（Tab 键）展开命令以及选项 \opt{-group} 与 \opt{-path}：
\begin{lstlisting}
pt_shell> report_timing -group SYS_CLK -path short
\end{lstlisting}

\begin{noteBox}
\begin{itemize}
  \item 数据到达（data arrival）段中包含数据通路单元及其延迟的行被去掉，因此报告称为“short”。
  \item 上述命令默认生成 setup 报告。
\end{itemize}
\end{noteBox}

\begin{questionBox}
\textbf{问题 20.}\ 至少有 4 处线索表明这是 setup 而非 hold 报告。你能找出多少？\\[0.35em]
\textbf{答：} 最醒目的线索是报告中的 \textbf{library setup time}。更细微的线索包括：
\begin{itemize}
  \item 数据到达与数据要求所用的时钟边分别为 0\,ns 与 8\,ns（而非 hold 时的 0\,ns 与 0\,ns）；
  \item 报告头中的 Path Type 为 \textbf{max}，表示 setup（\textbf{min} 表示 hold）；
  \item 数据到达时间晚于数据要求时间，且 slack 违例（若为 hold，在类似情形下 slack 往往表现为不同关系）。
\end{itemize}

\textbf{问题 21.}\ 写出起点与终点触发器的实例名。\\[0.35em]
\textbf{答：}\\
\begin{lstlisting}
Startpoint: I_ORCA_TOP/I_BLENDER/s3_op2_reg[18]
  (rising edge-triggered flip-flop clocked by SYS_CLK)
Endpoint:   I_ORCA_TOP/I_BLENDER/s4_op2_reg[31]
  (rising edge-triggered flip-flop clocked by SYS_CLK)
Path Group: SYS_CLK
Path Type:  max
\end{lstlisting}

\textbf{问题 22.}\ 该时序路径的时钟偏斜（clock skew）为 0.511\,ns；可用报告中的哪两行计算得到？\\[0.35em]
\textbf{答：} 可用启动端与捕获端的 \texttt{clock network delay (propagated)} 两行：
$3.247 - 2.736 = 0.511$\,ns。
\end{questionBox}

\begin{questionBox}
\textbf{问题 23.}\ 该时钟偏斜如何影响松弛（slack）——有助于还是有害于 slack？\\[0.35em]
\textbf{答：} 在本条 setup 报告中，时钟偏斜\textbf{有害于} slack：
启动端触发器的时钟延迟使数据晚到约 0.511\,ns，从而加大 setup 违例。

\textbf{问题 24.}\ 违例相对于时钟周期有多大？\\[0.35em]
\textbf{答：} 时钟周期为 8\,ns，违例约 0.987\,ns，大约相当于周期的 \textbf{12\%}。
\end{questionBox}

  \item 生成 hold 时序报告。\\
        下列快捷方式会执行历史中以字母 \texttt{rep} 开头的最近一条命令，并追加开关 \opt{-delay min}（从而生成 hold 报告）：
\begin{lstlisting}
pt_shell> !rep -delay min
\end{lstlisting}

\begin{questionBox}
\textbf{问题 25.}\ 至少有 4 处线索表明这是 hold 而非 setup 报告。你能找出多少？\\[0.35em]
\textbf{答：} 最醒目的线索是报告中的 \textbf{library hold time}。更细微的线索包括：
\begin{itemize}
  \item 数据到达与数据要求所用时钟边均为 0\,ns；
  \item Path Type 为 \textbf{min}，表示 hold；
  \item 数据到达早于数据要求时间，且 slack 违例。
\end{itemize}

\textbf{问题 26.}\ 在本 hold 报告中，时钟偏斜如何影响 slack——有助于还是有害于 slack？\\[0.35em]
\textbf{答：} 从报告可见（施加 clock network delay 之后），捕获时钟边晚于启动时钟边到达，
这\textbf{有害于} hold slack，并导致违例。
\end{questionBox}

  \item 退出 PrimeTime：
\begin{lstlisting}
pt_shell> quit
\end{lstlisting}
\end{enumerate}

\noindent 至此 Lab~1 完成。请返回课堂讲义。
